\documentclass[11pt,a4paper]{beamer}
\usetheme{Madrid}
\usepackage[utf8]{inputenc}
\usepackage{amsmath}
\usepackage{amsfonts}
\usepackage{amssymb}
\usepackage{graphicx}
\usepackage{xcolor}
\usepackage{tikz}
\usepackage{booktabs}
\usepackage{array}
\usepackage{colortbl}
\usepackage{multicol}


% ============================================================
% EXTREMELY SMALL FONT - MAXIMUM CONTENT PER SLIDE
% ============================================================
\usefonttheme{professionalfonts}
\setbeamerfont{itemize/enumerate body}{size=\tiny}
\setbeamerfont{itemize/enumerate subbody}{size=\tiny}
\setbeamerfont{block body}{size=\tiny}
\setbeamerfont{block title}{size=\scriptsize}
\setbeamerfont{frametitle}{size=\small}
\setbeamerfont{title}{size=\large}
\setbeamerfont{subtitle}{size=\small}
\setbeamerfont{author}{size=\tiny}
\setbeamerfont{date}{size=\tiny}

\renewcommand{\tiny}{\fontsize{5pt}{6pt}\selectfont}
\renewcommand{\scriptsize}{\fontsize{6pt}{7pt}\selectfont}
\renewcommand{\footnotesize}{\fontsize{6.5pt}{8pt}\selectfont}
\renewcommand{\small}{\fontsize{7.5pt}{9pt}\selectfont}
\renewcommand{\normalsize}{\fontsize{8.5pt}{10.5pt}\selectfont}

\newcommand{\examfocus}[1]{\textcolor{orange!80!black}{\textbf{[FOCUS] #1}}}
\newcommand{\trap}[1]{\textcolor{purple}{\textbf{[TRAP] #1}}}
\newcommand{\correct}[1]{\textcolor{green!50!black}{\textbf{[CORRECT] #1}}}
\newcommand{\scenario}[1]{\textcolor{blue!70!black}{\textbf{[SCENARIO] #1}}}
\newcommand{\keyterm}[1]{\textcolor{red!80!black}{\textbf{[KEY] #1}}}
\newcommand{\lawcite}[1]{\textcolor{brown!70!black}{\textbf{[LAW] #1}}}

\newcommand{\redflag}[1]{\textcolor{red!80!black}{\textbf{[RED FLAG: #1]}}}

\newcommand{\bluenote}[1]{\textcolor{blue!80!black}{\textbf{[NOTE: #1]}}}  % <-- ADD THIS LINE



\title[CAMS Domain C]{CAMS Exam Preparation\\Part C: Building an Anti-Financial Crime\\Compliance Program (30 percent of Exam)}
\author{Advanced AML Training Framework}
\date{}

\setbeamercovered{transparent}
\setbeamertemplate{navigation symbols}{}
\setbeamertemplate{frametitle continuation}{}
\setlength{\parskip}{2pt}
\setlength{\itemsep}{1pt}

\begin{document}

% ============================================================
% TITLE SLIDE
% ============================================================
\begin{frame}
	\titlepage
	\centering
	\tiny Domain 6: 30 percent of CAMS Exam | 110+ Slides | Governance, Controls, Investigations and Testing
\end{frame}

% ============================================================
% SECTION 1: THREE LINES OF DEFENSE (15 SLIDES)
% ============================================================

\begin{frame}{6.1: Three Lines of Defense - Complete Framework}
	\begin{block}{Core Concept}
		A foundational governance model ensuring separation between risk-taking, risk oversight, and independent assurance.
	\end{block}
	\begin{tabular}{p{0.18\textwidth}|p{0.32\textwidth}|p{0.45\textwidth}}
		\hline
		\rowcolor{lightgray}
		\textbf{Line} & \textbf{Who} & \textbf{Primary Role} \\
		\hline
		\textbf{First Line} & Business ops, RMs, tellers, loan officers & Own and manage risk; identify anomalies; complete initial CDD; escalate to compliance \\
		\hline
		\textbf{Second Line} & Compliance Dept, BSA Officer, MLRO & Oversee risk; set rules; monitor First Line; investigate; file SARs \\
		\hline
		\textbf{Third Line} & Internal Audit, external auditors & Independent assurance; audit First and Second Lines; report to Board \\
		\hline
	\end{tabular}
	\vspace{0.2cm}
	\redflag{Golden Rule: No single line should both DESIGN and TEST the same controls.}
	\examfocus{Exam tests whether you can identify independence violations across all three lines.}
\end{frame}

\begin{frame}{6.1.1: First Line - Detailed Responsibilities}
	\begin{block}{First Line (Business Operations) - What They MUST Do}
		The First Line is the front door of AML defense. They interact with customers directly.
	\end{block}
	\begin{multicols}{2}
		\textbf{Specific Duties:}
		\begin{itemize}
			\item Complete initial CIP
			\item Assess expected activity profile
			\item Identify unusual transactions
			\item Escalate concerns to compliance
			\item Respond to compliance RFIs
			\item Maintain customer records
			\item Update changes in customer status
		\end{itemize}
		\columnbreak
		\textbf{Common Red Flags to Spot:}
		\begin{itemize}
			\item Reluctance to provide ID
			\item Unusual transaction patterns
			\item Structuring cash deposits
			\item Changes without explanation
			\item Inconsistent narratives
			\item Third-party transactions
			\item Reluctance to explain source
		\end{itemize}
	\end{multicols}
	\redflag{Critical: First Line cannot close alerts or make SAR decisions. That is Second Line authority.}
\end{frame}

\begin{frame}{6.1.2: First Line Violation - RM Ignores Red Flags}
	\scenario{RM Tells Analyst to Close Alert Based on Personal Knowledge}
	A relationship manager has known a corporate client for 12 years. The client suddenly makes weekly cash deposits of \$9,500 - just below the \$10,000 CTR threshold. The RM tells compliance: I have known this customer for over a decade. He is trustworthy. Just close the alert. The analyst closes the alert without independent review, documentation, or supervisor consultation.
	
	\vspace{0.1cm}
	\textbf{How to Analyze:}
	\begin{enumerate}
		\item First Line's job: IDENTIFY and ESCALATE - not investigate or close
		\item Second Line's job: INVESTIGATE independently
		\item The deposit pattern (\$9,500 weekly) is structuring
		\item Closing without documentation violates AML program requirements
	\end{enumerate}
	\trap{RM's long relationship is NOT sufficient evidence to close an alert.}
	\correct{Compliance must conduct independent investigation with documented evidence.}
\end{frame}

\begin{frame}{6.1.3: Second Line - Detailed Responsibilities}
	\begin{block}{Second Line (Compliance) - Core Duties}
		The Second Line must operate independently from the First Line and have authority to override business decisions.
	\end{block}
	\begin{multicols}{2}
		\textbf{Core Duties:}
		\begin{itemize}
			\item Design AML policies
			\item Conduct transaction monitoring
			\item Investigate alerts
			\item Make SAR filing decisions
			\item File SARs with FIU
			\item Provide training
			\item Conduct EDD
			\item Maintain risk assessment
			\item Respond to law enforcement
			\item Report to Board
		\end{itemize}
		\columnbreak
		\textbf{Independence Requirements:}
		\begin{itemize}
			\item Cannot report to business lines
			\item Direct Board access
			\item Override business decisions
			\item Sufficient staffing
			\item No commercial pressure
			\item Access to all information
			\item Independent budget
		\end{itemize}
	\end{multicols}
	\redflag{Commercial importance of a client is NEVER a valid reason to delay or avoid SAR filing.}
\end{frame}

\begin{frame}{6.1.4: Second Line Failure - No Documentation}
	\scenario{Analyst Closes Alert Without Investigation}
	An analyst reviews an alert for a restaurant customer with 23 cash deposits between $8,400 and $9,800 over 45 days. The analyst writes: Restaurant business. Cash is normal. No SAR needed. The analyst closes without: daily sales logs, prior period comparison, supervisor consultation, or any documentation.
	
	\vspace{0.1cm}
	\textbf{How to Analyze:}
	\begin{enumerate}
		\item 23 deposits just below \$10,000 = structuring indicator
		\item Structuring is a federal crime regardless of funds origin
		\item The analyst assumed legitimacy without verification
		\item No documentation = no investigation occurred
	\end{enumerate}
	\trap{Analyst used reasonable business judgment based on industry.}
	\correct{Analysts must obtain supporting documentation and document all investigative steps. Assumptions are not investigations.}
\end{frame}

\begin{frame}{6.1.5: Third Line - Independence Requirements}
	\begin{block}{Internal Audit - What They CAN and CANNOT Do}
		Audit independence is absolutely critical to program credibility.
	\end{block}
	\begin{tabular}{p{0.45\textwidth}|p{0.45\textwidth}}
		\hline
		\rowcolor{lightgray}
		\textbf{Audit CAN Do} & \textbf{Audit CANNOT Do} \\
		\hline
		Test effectiveness of controls & Design or build controls \\
		Recommend improvements & Operate monitoring systems \\
		Report findings to Board & Make SAR decisions \\
		Review sample of alerts & Approve customer relationships \\
		Assess policy compliance & Perform ongoing monitoring \\
		Evaluate training & Conduct EDD on customers \\
		\hline
	\end{tabular}
	\redflag{Most Common Violation: Audit helping to design scenarios THEN auditing those same scenarios.}
\end{frame}

\begin{frame}{6.1.6: Third Line Independence Violation}
	\scenario{Audit Designs Controls Then Tests Them}
	Compliance lacks technical expertise. CCO asks Internal Audit for help. Audit spends 3 months co-designing 14 new monitoring scenarios. Six months later, the same audit team tests the system and finds all 14 scenarios effective.
	
	\vspace{0.1cm}
	\textbf{How to Analyze:}
	\begin{enumerate}
		\item Audit's role is to TEST, not DESIGN
		\item Designing controls makes Audit part of Second Line
		\item Auditors cannot objectively test their own work
		\item The audit findings are unreliable
	\end{enumerate}
	\trap{Scenarios were effective, so no violation. OR CCO appropriately sought help.}
	\correct{Audit must remain independent. Use external consultants or different internal team with no design involvement.}
\end{frame}

\begin{frame}{6.1.7: Board of Directors - What They MUST Do}
	\begin{block}{Board Oversight - Ultimate Accountability}
		The Board is ultimately accountable for the AML program. Regulators hold the Board responsible.
	\end{block}
	\begin{multicols}{2}
		\textbf{Board MUST:}
		\begin{itemize}
			\item Approve initial AML program
			\item Approve material changes
			\item Appoint BSA Officer
			\item Receive regular updates
			\item Be informed of deficiencies
			\item Hold management accountable
			\item Review audit findings
			\item Ensure adequate resources
		\end{itemize}
		\columnbreak
		\textbf{Board DOES NOT:}
		\begin{itemize}
			\item File SARs
			\item Review KYC files
			\item Monitor daily alerts
			\item Conduct investigations
			\item Make SAR decisions
			\item Perform testing
			\item Approve transactions
		\end{itemize}
	\end{multicols}
	\redflag{Board approval is required for the AML program. CCO cannot approve their own program.}
\end{frame}

\begin{frame}{6.1.8: Board Violation - Hiding Deficiencies}
	\scenario{Management Withholds Information from Board}
	Compliance discovers a 14-month monitoring gap affecting \$340 million in wires. Senior management decides not to inform the Board because the issue is operational and will be fixed within 90 days.
	
	\vspace{0.1cm}
	\textbf{How to Analyze:}
	\begin{enumerate}
		\item 14-month gap affecting \$340 million is MATERIAL
		\item Material deficiencies MUST be reported to Board immediately
		\item Management cannot filter or delay reporting
		\item Board needs this information for oversight
	\end{enumerate}
	\trap{Senior management has authority to decide what reaches the Board.}
	\correct{Board must be informed of material deficiencies. Management cannot withhold this information.}
\end{frame}

\begin{frame}{6.1.9: Board vs. Management - Approval Authority}
	\scenario{CCO Implements Policy Without Board Approval}
	CCO drafts revised AML policy. Legal reviews and approves. CCO implements without presenting to Board. CCO argues: I am the designated BSA Officer with authority to implement policies. Board approval is unnecessary.
	
	\vspace{0.1cm}
	\textbf{How to Analyze:}
	\begin{enumerate}
		\item Board must approve AML program under BSA
		\item Material changes require Board approval
		\item CCO implements, does not have final approval
		\item Policy is not effective without Board approval
	\end{enumerate}
	\trap{CCO has delegated authority to implement policies.}
	\correct{Board must approve AML program and material changes. CCO implements Board-approved program but lacks final approval authority.}
\end{frame}

\begin{frame}{6.1.10: Senior Management Responsibilities}
	\begin{block}{Senior Management - Execution and Accountability}
		Senior management executes the Board-approved program and is accountable to the Board.
	\end{block}
	\begin{multicols}{2}
		\textbf{Senior Management MUST:}
		\begin{itemize}
			\item Implement Board-approved program
			\item Allocate sufficient resources
			\item Ensure timely SAR filings
			\item Respond to audit findings
			\item Report material issues to Board
			\item Appoint qualified staff
			\item Ensure training effectiveness
		\end{itemize}
		\columnbreak
		\textbf{Senior Management CANNOT:}
		\begin{itemize}
			\item Veto SARs for commercial reasons
			\item Withhold information from Board
			\item Retaliate against whistleblowers
			\item Operate without Board approval
			\item Delegate accountability
			\item Ignore audit findings
		\end{itemize}
	\end{multicols}
	\redflag{Senior management faces personal regulatory risk for AML failures, including fines and imprisonment.}
\end{frame}

% ============================================================
% SECTION 2: INDEPENDENT TESTING (10 SLIDES)
% ============================================================

\begin{frame}{6.2: Independent Testing - What Must Be Tested}
	\begin{block}{Required Scope - Operational Effectiveness}
		Independent testing is one of the five mandatory BSA pillars. Must test OPERATIONAL effectiveness.
	\end{block}
	\begin{tabular}{p{0.3\textwidth}|p{0.65\textwidth}}
		\hline
		\rowcolor{lightgray}
		\textbf{Area} & \textbf{Specific Testing Procedures} \\
		\hline
		Transaction Monitoring & Test thresholds using historical SARs; calculate false negative rate; validate alert logic \\
		\hline
		Sanctions Screening & Test fuzzy logic; validate name matching; review override decisions \\
		\hline
		CDD/KYC & Review high-risk samples; verify UBO identification; test EDD documentation \\
		\hline
		SAR Filing & Review timeliness; assess quality; test escalation; validate no tipping-off \\
		\hline
		Training & Assess role-based content; test completion rates; evaluate effectiveness \\
		\hline
		Governance & Review Board approvals; test independence; review escalation documentation \\
		\hline
	\end{tabular}
	\redflag{Not sufficient: Only reviewing policies without testing transactions.}
\end{frame}

\begin{frame}{6.2.1: Testing Transaction Monitoring - False Negatives}
	\scenario{Audit Discovers 77 Percent False Negative Rate}
	Auditor tests system by running 500 historical SAR cases through current system. 387 of 500 (77 percent) would NOT generate an alert under current thresholds. Compliance director argues this is acceptable because bank already reviews 22,000 alerts per month.
	
	\vspace{0.1cm}
	\textbf{How to Analyze:}
	\begin{enumerate}
		\item 77 percent false negative rate is EXTREMELY high
		\item System would miss over three-quarters of confirmed suspicious activity
		\item Operational burden is NOT a valid excuse
		\item Regulators prioritize detection over efficiency
	\end{enumerate}
	\trap{High alert volume justifies maintaining current thresholds.}
	\correct{System must be effective first. Document deficiency, create remediation plan, recalibrate thresholds.}
\end{frame}

\begin{frame}{6.2.2: Who Can Perform Independent Testing?}
	\begin{block}{Qualification Requirements}
		The tester must meet specific independence and expertise requirements.
	\end{block}
	\begin{tabular}{p{0.45\textwidth}|p{0.45\textwidth}}
		\hline
		\rowcolor{lightgray}
		\textbf{Acceptable} & \textbf{NOT Acceptable} \\
		\hline
		Internal Audit (if independent) & Compliance testing itself \\
		External Audit Firm & Audit that designed controls \\
		Third-Party Consultant & Former compliance staff who designed policies \\
		Independent consultant & Business line employees \\
		No prior involvement & Vendors who built the system \\
		\hline
	\end{tabular}
	\textbf{Requirements:}
	\begin{itemize}
		\item Independence: No involvement in designing or operating controls
		\item Expertise: AML/CFT knowledge and qualifications
		\item Objectivity: No conflicts of interest
		\item Reporting: Reports to Board or Audit Committee, not CCO
	\end{itemize}
\end{frame}

\begin{frame}{6.2.3: Prior Involvement Trap}
	\scenario{Tester Previously Worked in Compliance}
	Bank hires external consulting firm for AML audit. Firm assigns individual who was bank's BSA Officer until 8 months ago. This individual personally designed several policies now being audited. CCO sees no issue because consultant is external.
	
	\vspace{0.1cm}
	\textbf{How to Analyze:}
	\begin{enumerate}
		\item Independence is about the PERSON, not just the firm
		\item Person who designed controls cannot audit those controls
		\item 8-month separation does NOT restore objectivity
		\item Firm should assign different personnel
	\end{enumerate}
	\trap{Consultant is external, so independence is automatically satisfied.}
	\correct{Individual tester must have no prior involvement. Assign different team member.}
\end{frame}

\begin{frame}{6.2.4: Audit Reporting Line - CCO Cannot Approve Findings}
	\scenario{Audit Reports to CCO for Approval}
	Audit completes annual AML audit and issues draft report identifying deficiencies. Audit manager provides draft to CCO for approval before finalizing. CCO requests that critical findings be removed. Audit manager agrees. Final report shows no deficiencies.
	
	\vspace{0.1cm}
	\textbf{How to Analyze:}
	\begin{enumerate}
		\item Audit must report findings directly to Board
		\item Management cannot approve or edit audit findings
		\item Removing findings destroys audit independence
		\item Final report is unreliable
	\end{enumerate}
	\trap{CCO is responsible for AML program and should review findings for accuracy.}
	\correct{Audit findings must be reported directly to Board without management filtering. CCO may receive copy but cannot control content.}
\end{frame}

\begin{frame}{6.2.5: Audit Frequency and Risk-Based Approach}
	\begin{block}{How Often Must Testing Be Performed?}
		Frequency depends on risk profile, but at least annually for most institutions.
	\end{block}
	\textbf{Risk-Based Frequency Considerations:}
	\begin{itemize}
		\item High-risk institutions: Semi-annual or quarterly testing recommended
		\item Medium-risk institutions: Annual testing is standard
		\item Low-risk institutions: Annual testing is minimum, but may be less intensive
		\item Material changes: Trigger interim testing outside normal cycle
		\item Regulatory findings: May require more frequent testing
	\end{itemize}
	\textbf{Trigger Events for Interim Testing:}
	\begin{itemize}
		\item New high-risk products or services
		\item Significant customer base growth
		\item Merger or acquisition
		\item Regulatory enforcement action
		\item Material system changes
		\item Emergence of new typologies
		\item Known control deficiencies
	\end{itemize}
	\redflag{Annual is minimum, not maximum. Material changes require interim testing.}
\end{frame}

% ============================================================
% SECTION 3: TRANSACTION MONITORING CALIBRATION (12 SLIDES)
% ============================================================

\begin{frame}{6.3: False Positives vs. False Negatives}
	\begin{block}{The Core Calibration Trade-Off}
		Every AML monitoring system faces this fundamental decision.
	\end{block}
	\begin{tabular}{p{0.25\textwidth}|p{0.35\textwidth}|p{0.35\textwidth}}
		\hline
		\rowcolor{lightgray}
		\textbf{Term} & \textbf{Definition} & \textbf{Risk} \\
		\hline
		False Positive & Legitimate activity triggers alert & Wasted analyst time; operational inefficiency; analyst fatigue \\
		\hline
		False Negative & Suspicious activity generates no alert & MISSED criminal activity; direct AML deficiency; enforcement risk \\
		\hline
	\end{tabular}
	\vspace{0.2cm}
	\redflag{Which is worse? FALSE NEGATIVES are significantly worse.}
	\bluenote{Regulator Priority: Detection effectiveness > operational efficiency.}
\end{frame}

\begin{frame}{6.3.1: Over-Filtering Trap - Raise Thresholds Too High}
	\scenario{Director Raises Thresholds by 250 Percent}
	Bank generates 8,000 alerts/month with 96 percent false positives. Director orders all thresholds raised by 250 percent to reduce costs. Team leader warns this will increase false negatives. Director replies: Any missed activity is acceptable as long as we still file some SARs.
	
	\vspace{0.1cm}
	\textbf{How to Analyze:}
	\begin{enumerate}
		\item 250 percent increase is extreme and arbitrary
		\item Many suspicious transactions will fall below thresholds
		\item Director prioritizing cost over effectiveness
		\item Regulators will view as material deficiency
	\end{enumerate}
	\trap{Director correctly managing operational efficiency under risk-based approach.}
	\correct{Raising thresholds must be data-driven. Validate new thresholds still detect known typologies. Effectiveness cannot be sacrificed.}
\end{frame}

\begin{frame}{6.3.2: Under-Filtering Trap - Too Many Alerts}
	\scenario{System Generates 50,000 Alerts Per Month}
	Bank sets thresholds very low to catch everything. System generates 50,000 alerts/month. Team has 5 analysts, each can review 200 alerts/month. Backlog grows to 40,000 uninvestigated alerts. Director argues lower thresholds are always better.
	
	\vspace{0.1cm}
	\textbf{How to Analyze:}
	\begin{enumerate}
		\item Team can only review 1,000 alerts/month
		\item 49,000 alerts/month are never reviewed
		\item Unreviewed alerts provide zero detection value
		\item Bank must increase staffing OR recalibrate
	\end{enumerate}
	\trap{Lower thresholds always better because they catch more activity.}
	\correct{System must be calibrated so team can review all alerts. Unreviewed alerts are a control failure.}
\end{frame}

\begin{frame}{6.3.3: Model Drift - Performance Degradation}
	\scenario{Model Performance Degrades Over 3.5 Years}
	Bank deployed ML model January 2022 with 72 percent precision and 81 percent recall. June 2025 validation shows 51 percent precision and 62 percent recall. Bank introduced new products, volumes up 180 percent, typologies shifted. Model never recalibrated.
	
	\vspace{0.1cm}
	\textbf{How to Analyze:}
	\begin{enumerate}
		\item Models degrade when behavior/products/typologies change
		\item 3.5 years without recalibration is a significant gap
		\item Annual recalibration or trigger-based retraining required
		\item Performance degradation requires immediate remediation
	\end{enumerate}
	\trap{Model still has recall above 60 percent, which is acceptable.}
	\correct{Bank must establish model monitoring and refresh processes. Significant degradation requires recalibration or replacement.}
\end{frame}

\begin{frame}{CAMS-Style Question: Transaction Monitoring System Effectiveness}
	\scenario{Bank Implemented System in 2022, No Updates by Mid-2025}
	Bank implemented automated transaction monitoring system in January 2022. By mid-2025, bank introduced new products, transaction volumes increased significantly, and financial crime typologies evolved. Recent internal review found system generating less effective results. Bank has not updated or reviewed the system since implementation.
	
	\vspace{0.1cm}
	\textbf{How to Analyze:}
	\begin{enumerate}
		\item Systems degrade when products/volumes/typologies change
		\item 3.5 years without review is a significant compliance gap
		\item Periodic tuning and updates are AML program requirements
		\item "Acceptable" performance without adaptation increases risk exposure
	\end{enumerate}
	\trap{Continue using the system since it is still identifying suspicious activity at an acceptable level.}
	\correct{Establish a process to periodically review, tune, and update the monitoring system.}
\end{frame}

\begin{frame}{6.3.4: Trigger-Based Validation Requirements}
	\begin{block}{When Must Interim Validation Be Performed?}
		Annual validation is minimum. Certain events trigger interim validation.
	\end{block}
	\textbf{Trigger Events Requiring Interim Validation:}
	\begin{itemize}
		\item New high-risk products or services (crypto, correspondent banking, private banking)
		\item Significant customer base growth (acquisition, organic growth over 25 percent)
		\item Emerging typologies (sudden increase in specific SAR types)
		\item Regulatory changes (new requirements or guidance)
		\item Model performance alerts (vendor notification of degradation)
		\item Material system changes (new vendor, major upgrade)
		\item Known control failures (audit findings, regulatory exam findings)
		\item Geographic expansion (new high-risk jurisdictions)
	\end{itemize}
	\textbf{Documentation Requirements:}
	\begin{itemize}
		\item Date of trigger event
		\item Validation methodology used
		\item Results of validation testing
		\item Any threshold adjustments made
		\item Approval by compliance management
		\item Board notification for material changes
	\end{itemize}
	\redflag{Policies that only require annual validation are deficient without trigger provisions.}
\end{frame}

\begin{frame}{6.3.5: Scenario Effectiveness Testing Methods}
	\begin{block}{How to Test If Scenarios Are Working}
		Multiple methods are used to validate scenario effectiveness.
	\end{block}
	\textbf{Historical SAR Testing (Back-testing):}
	\begin{itemize}
		\item Run historical SAR cases through current scenarios
		\item Calculate capture rate (percentage that would generate alerts)
		\item Acceptable capture rate depends on typology but should exceed 70 percent
		\item Low capture rate indicates over-filtering or scenario decay
	\end{itemize}
	\textbf{Holdout Sample Testing:}
	\begin{itemize}
		\item Reserve portion of SAR data not used in development
		\item Test scenarios on holdout sample to validate performance
		\item Helps detect overfitting to development data
	\end{itemize}
	\textbf{Scenario and Threshold Tuning:}
	\begin{itemize}
		\item Adjust parameters based on testing results
		\item Document all changes with rationale
		\item Re-test after changes to validate improvement
		\item Maintain version control of scenarios
	\end{itemize}
	\textbf{Benchmarking:}
	\begin{itemize}
		\item Compare detection rates to industry peers (if data available)
		\item Use regulatory guidance on expected detection rates
		\item Consider typology-specific capture rates
	\end{itemize}
\end{frame}

% ============================================================
% SECTION 4: DATA QUALITY AND RECONCILIATION (8 SLIDES)
% ============================================================

\begin{frame}{6.4: Data Quality - The Silent Failure}
	\begin{block}{Core Concept}
		A monitoring system is only as good as the data it receives. Missing data = missing alerts.
	\end{block}
	\textbf{Critical Controls That Must Exist:}
	\begin{itemize}
		\item Reconciliation between source systems and monitoring system (daily or real-time)
		\item Data completeness and accuracy checks (automated validation rules)
		\item Error handling and alerting for failed data feeds (no silent failures)
		\item Regular validation that all required transactions are ingested (sample testing)
		\item Data lineage documentation (know where each data element comes from)
		\item Data quality metrics and reporting to management
	\end{itemize}
	\redflag{The worst data failure: Silent failure where data is missing but no error messages alert the team.}
	\examfocus{Exam tests whether you understand that data quality is a compliance responsibility, not just IT.}
\end{frame}

\begin{frame}{6.4.1: Silent Data Feed Failure - Complete Analysis}
	\scenario{34 Percent of Wires Missing for 14 Months}
	Bank's monitoring system relies on daily feeds from 17 upstream systems. Audit reveals that for 14 months, the feed from the international wire system silently failed for 34 percent of all wires. Those transactions never entered monitoring, never generated alerts, never reviewed. Estimated \$340 million in wires not monitored. Failure was technical misconfiguration with no error messages. No reconciliation controls existed to detect the gap.
	
	\vspace{0.1cm}
	\textbf{How to Analyze:}
	\begin{enumerate}
		\item 34 percent missing data = massive monitoring gap
		\item No alerts for transactions system never received
		\item Absence of reconciliation controls is foundational failure
		\item Bank must remediate AND review missed transactions retrospectively
		\item SARs must be filed for any suspicious activity identified in missed transactions
		\item Regulators will view this as material deficiency
	\end{enumerate}
	\trap{Misconfiguration was accidental, so no compliance violation occurred.}
	\correct{Bank must have controls to verify complete data ingestion. 14-month gap requires retrospective review and SAR filing.}
\end{frame}

\begin{frame}{6.4.2: Data Reconciliation - Required Controls}
	\begin{block}{Reconciliation Controls - What Must Exist}
		Banks must be able to prove that all transactions entered the monitoring system.
	\end{block}
	\textbf{Reconciliation Methodology:}
	\begin{itemize}
		\item Source system transaction count vs. monitoring system intake count (daily)
		\item Source system total dollar volume vs. monitoring system total volume (daily)
		\item Hash total comparisons for data integrity
		\item Exception report for any discrepancies
		\item Timely resolution of exceptions (24-48 hours)
		\item Escalation for unresolved exceptions
	\end{itemize}
	\textbf{Common Data Quality Issues:}
	\begin{itemize}
		\item Missing transactions (data feed failure)
		\item Duplicate transactions (processing errors)
		\item Corrupted data (format issues)
		\item Delayed data (latency beyond acceptable window)
		\item Incomplete data (missing required fields)
		\item Incorrect data (wrong values, mapping errors)
	\end{itemize}
	\textbf{Risk of No Reconciliation:}
	\begin{itemize}
		\item Cannot demonstrate monitoring completeness to regulators
		\item Undetected data gaps may exist for months or years
		\item Criminal activity may be processed without detection
		\item Bank may face enforcement action for monitoring failure
	\end{itemize}
\end{frame}

\begin{frame}{6.4.3: Missing Customer Identification - Unlinked Transactions}
	\scenario{12 Percent of Transactions Have NULL Customer ID}
	Monitoring system flags large wire transfers. When analyst investigates, system shows Customer ID: NULL value. Data issue present for 8 months. IT explains that mobile app account openings sometimes fail to populate unique identifier. Approximately 12 percent of transactions cannot be linked to any customer record.
	
	\vspace{0.1cm}
	\textbf{How to Analyze:}
	\begin{enumerate}
		\item Unlinked transactions cannot be investigated properly
		\item Bank cannot determine customer risk or conduct CDD
		\item Data quality failure with direct AML impact
		\item Cannot file accurate SARs without customer identification
		\item Root cause must be fixed
		\item Process needed to identify and investigate orphaned transactions
	\end{enumerate}
	\trap{IT issues are operational and not a compliance concern.}
	\correct{Missing customer identification is material AML deficiency. Fix root cause and develop process for orphaned transactions.}
\end{frame}

\begin{frame}{6.4.4: Data Quality Metrics and Governance}
	\begin{block}{Data Quality Key Performance Indicators}
		Compliance must monitor and report data quality metrics to management.
	\end{block}
	\textbf{Essential Data Quality Metrics:}
	\begin{itemize}
		\item Data completeness percentage by source system (target: 99.9 percent+)
		\item Data timeliness (average latency from transaction to monitoring)
		\item Data accuracy (error rate in critical fields)
		\item Reconciliation exception rate and aging
		\item Number of unlinked transactions
		\item Number of failed data feeds
	\end{itemize}
	\textbf{Governance Requirements:}
	\begin{itemize}
		\item Data quality is shared responsibility (IT, Compliance, Business)
		\item Compliance must define data requirements for monitoring
		\item IT must implement controls to meet requirements
		\item Business must ensure source system data quality
		\item Regular data quality review meetings
		\item Escalation path for unresolved data issues
		\item Board reporting on material data gaps
	\end{itemize}
	\redflag{Compliance cannot simply accept whatever data IT provides. Compliance must define requirements.}
\end{frame}

% ============================================================
% SECTION 5: ALERT ADJUDICATION (6 SLIDES)
% ============================================================

\begin{frame}{6.5: Alert Adjudication - Investigator Variation Problem}
	\begin{block}{Core Concept}
		Different analysts should reach similar conclusions on similar facts. Wide variation indicates lack of standards.
	\end{block}
	\textbf{Required Controls to Ensure Consistency:}
	\begin{itemize}
		\item Written adjudication guidelines for each scenario
		\item Required documentation standards (what must be in every case file)
		\item Quality assurance reviews of closed alerts (sample testing)
		\item Calibration sessions to align analyst judgment
		\item Supervisor approval for certain dispositions
		\item Appeal process for disputed decisions
		\item Regular training on adjudication standards
	\end{itemize}
	\redflag{Red Flag: One analyst closes 95 percent of alerts as false positives; another escalates 70 percent for the same scenario.}
	\examfocus{Exam tests whether you know that standardization is required, not optional.}
\end{frame}

\begin{frame}{6.5.1: Inconsistent Analyst Problem - Complete Analysis}
	\scenario{94 Percent vs. 67 Percent Escalation Rate}
	QA team reviews 200 closed alerts for cash deposit scenario. Analyst A closed 94 percent as false positives with note: Customer is a retailer; cash is normal. Analyst B, reviewing similar accounts for same scenario, escalated 67 percent for further investigation and filed SARs on 22 percent. Neither analyst's file includes supporting documentation. Bank has no written adjudication guidelines for this scenario.
	
	\vspace{0.1cm}
	\textbf{How to Analyze:}
	\begin{enumerate}
		\item 94 percent vs. 67 percent variation is extreme and unacceptable
		\item Without written guidelines, analysts apply inconsistent standards
		\item No documentation = bank cannot demonstrate investigations occurred
		\item Bank needs written guidelines AND QA oversight
		\item Analyst A may be under-escalating (missing suspicious activity)
		\item Analyst B may be over-escalating (wasting resources)
		\item Either way, process is not controlled
	\end{enumerate}
	\trap{Different analysts apply professional judgment differently; variation is acceptable.}
	\correct{Bank must have written adjudication standards, required documentation, and QA processes.}
\end{frame}

\begin{frame}{6.5.2: Uninvestigated Alert - Assuming Without Evidence}
	\scenario{Analyst Closes Alert Based on Assumption}
	Alert flags customer with 23 cash deposits between $8,400 and $9,700 over 45 days. Customer is restaurant owner. Analyst writes: Restaurant business. Cash from customers. No SAR needed. Analyst closes without requesting daily sales logs, comparing to prior periods, checking if pattern changed after new manager hired, or any supporting documentation.
	
	\vspace{0.1cm}
	\textbf{How to Analyze:}
	\begin{enumerate}
		\item Deposit pattern is classic structuring indicator
		\item Analyst assumed legitimacy without verification
		\item No supporting documentation was obtained
		\item Alert closed without meaningful investigation
		\item Cannot defend this decision in regulatory exam
	\end{enumerate}
	\trap{Analyst used business context to make reasonable judgment.}
	\correct{Analysts must obtain and review supporting documentation. Assumptions without verification do not constitute investigation.}
\end{frame}

\begin{frame}{6.5.3: Quality Assurance Program Requirements}
	\begin{block}{QA Program - Essential Elements}
		Quality assurance is the check on analyst decision-making.
	\end{block}
	\textbf{QA Program Components:}
	\begin{itemize}
		\item Sample size: Sufficient to be statistically valid (typically 5-10 percent of alerts)
		\item Sampling methodology: Random plus targeted (high-risk scenarios, specific analysts)
		\item QA reviewer independence: Different from original analyst
		\item QA reviewer qualifications: More experienced than front-line analysts
		\item Dispute resolution process: Path for disagreement between analyst and QA
		\item Remediation: Process for correcting errors found by QA
		\item Trending: Track QA findings over time to identify patterns
		\item Reporting: QA results reported to compliance management monthly
	\end{itemize}
	\textbf{What QA Reviews:}
	\begin{itemize}
		\item Did analyst obtain required documentation?
		\item Is documentation sufficient to support conclusion?
		\item Was analysis thorough and logical?
		\item Was conclusion reasonable based on evidence?
		\item Were escalation procedures followed?
		\item Was SAR filed when required?
		\item Was tipping-off avoided?
	\end{itemize}
	\redflag{No QA program = no assurance that analysts are making correct decisions.}
\end{frame}

% ============================================================
% SECTION 6: SAR FILING AND TIPPING-OFF (10 SLIDES)
% ============================================================

\begin{frame}{6.6: SAR Filing - Timing and Legal Standard}
	\begin{block}{Core Concept}
		SARs must be filed immediately upon suspicion, regardless of transaction value.
	\end{block}
	\textbf{Key Principles:}
	\begin{itemize}
		\item Legal standard = REASONABLE SUSPICION (not certainty, not probable cause, not proof)
		\item No minimum monetary threshold for money laundering suspicion
		\item Do NOT delay filing to gather more evidence
		\item Do NOT investigate to point of certainty before filing
		\item File first, continue investigating if needed
		\item Cannot consider commercial impact of filing
		\item Cannot notify customer (tipping-off prohibition)
	\end{itemize}
	\textbf{What Reasonable Suspicion Means:}
	\begin{itemize}
		\item More than a hunch, less than probable cause
		\item Specific articulable facts that suggest possible money laundering
		\item Based on training, experience, and available information
		\item Does not require knowledge of underlying predicate offense
		\item Does not require confirmation that funds are criminal
	\end{itemize}
	\redflag{Common violation: We need more evidence before filing. This is incorrect. Reasonable suspicion is enough.}
\end{frame}

\begin{frame}{6.6.1: Need More Evidence Trap - Complete Analysis}
	\scenario{Supervisor Delays SAR Filing for 60 Days}
	Compliance analyst identifies pattern raising reasonable suspicion of money laundering. Supervisor says: We do not have enough evidence yet. Keep investigating and file when you have conclusive proof. Analyst continues investigating for 60 days without filing SAR.
	
	\vspace{0.1cm}
	\textbf{How to Analyze:}
	\begin{enumerate}
		\item Reasonable suspicion is the legal standard, not proof
		\item Conclusive proof is NOT required for SAR filing
		\item Delaying filing is a regulatory violation
		\item Bank can file SAR and continue investigating
		\item Supervisor misstated the legal standard
		\item 60-day delay is significant and indefensible
	\end{enumerate}
	\trap{Supervisor correctly wants more evidence to ensure SAR is accurate.}
	\correct{SAR must be filed immediately upon reasonable suspicion. File first, then continue investigating. Delaying is violation.}
\end{frame}

\begin{frame}{6.6.2: Tipping-Off - What Constitutes a Violation}
	\scenario{Teller Tells Customer About Compliance Review}
	Customer asks teller why wire transfer is delayed. Teller says: I am sorry, but our compliance team is reviewing your account for some unusual activity. Customer becomes agitated and closes account. No SAR has been filed yet; investigation is ongoing.
	
	\vspace{0.1cm}
	\textbf{How to Analyze:}
	\begin{enumerate}
		\item Tipping-off includes disclosing that an investigation EXISTS
		\item You do NOT need to mention SAR by name to tip off
		\item Telling customer they are under compliance review IS tipping-off
		\item Violation occurred even though no SAR filed yet
		\item Teller should have given neutral response or said nothing
	\end{enumerate}
	\trap{No SAR was filed yet, so no tipping-off occurred.}
	\correct{Tipping-off includes disclosing existence of internal investigation. Train staff not to disclose compliance scrutiny.}
\end{frame}

\begin{frame}{6.6.3: Permissible Disclosures - 314(b) Sharing}
	\scenario{Two Banks Share SAR Information Under 314(b)}
	Two banks have filed SARs on different individuals who appear to be part of same money laundering network. Both banks have filed FinCEN 314(b) opt-in notifications. BSA Officers meet to share information about suspicious activity.
	
	\vspace{0.1cm}
	\textbf{How to Analyze:}
	\begin{enumerate}
		\item 314(b) permits voluntary information sharing between opted-in institutions
		\item Sharing is protected by safe harbor from liability
		\item This is NOT tipping-off because explicitly permitted by law
		\item Both banks must have filed opt-in notification
		\item Sharing must be for purpose of identifying ML/TF activity
	\end{enumerate}
	\trap{Sharing SAR information with another bank always constitutes tipping-off.}
	\correct{314(b) sharing is lawful exception to tipping-off. Banks that opted in can share information.}
\end{frame}

\begin{frame}{6.6.4: Internal Sharing - Legal Department Need-to-Know}
	\scenario{Analyst Shares SAR Draft with Legal Department}
	Compliance analyst drafts SAR. Before filing, analyst shares draft with bank's legal department for review. Legal department advises on potential legal risks. Analyst then files SAR.
	
	\vspace{0.1cm}
	\textbf{How to Analyze:}
	\begin{enumerate}
		\item Legal department has legitimate need to know
		\item Sharing is internal, not external
		\item SAR remains confidential, not disclosed to customer
		\item This is standard and permissible practice
		\item Does not constitute tipping-off
	\end{enumerate}
	\trap{Any sharing of SAR information before filing is tipping-off.}
	\correct{Internal sharing with legal, compliance management, and audit on need-to-know basis is permissible.}
\end{frame}

\begin{frame}{6.6.5: SAR Confidentiality - Criminal Penalties}
	\begin{block}{Unauthorized Disclosure is a Crime}
		Disclosing SAR information without authorization carries criminal penalties.
	\end{block}
	\textbf{Penalties for Unauthorized Disclosure:}
	\begin{itemize}
		\item Criminal fines up to \$250,000 for individuals
		\item Imprisonment up to 5 years
		\item Civil penalties for institutions
		\item Termination of employment
		\item Professional debarment
	\end{itemize}
	\textbf{Who Can Receive SAR Information:}
	\begin{itemize}
		\item FinCEN and other FIUs (by filing)
		\item Law enforcement (via proper channels)
		\item Other opted-in 314(b) institutions (limited)
		\item Legal department (internal need-to-know)
		\item Compliance management (internal need-to-know)
		\item Internal Audit (internal need-to-know)
		\item Board of Directors (oversight)
	\end{itemize}
	\textbf{Who CANNOT Receive SAR Information:}
	\begin{itemize}
		\item The customer (tipping-off)
		\item Customer's attorney (without legal process)
		\item Media or public
		\item Other banks not under 314(b)
		\item Foreign regulators without proper channel
	\end{itemize}
	\redflag{SAR confidentiality is absolute except for specific legal exceptions.}
\end{frame}

\begin{frame}{6.6.6: SAR Retention and Recordkeeping}
	\begin{block}{Recordkeeping Requirements for SARs}
		SARs and supporting documentation must be retained for specific periods.
	\end{block}
	\textbf{Retention Requirements:}
	\begin{itemize}
		\item SAR filing and all supporting documentation: 5 years from date of filing
		\item Copy of SAR (confidential file): 5 years

		\item RFIs and customer responses: 5 years
		\item Documentation of SAR decision: 5 years
	\end{itemize}
	\textbf{What Must Be Retained:}
	\begin{itemize}
		\item SAR form itself (if paper filed)
		\item Copy of electronic filing confirmation
		\item All transaction data reviewed
		\item Analyst investigation notes
		\item Supporting documentation obtained
		\item Supervisor approval documentation
		\item RFIs and responses
		\item Decision to file or not file
		\item Law enforcement feedback (if any)
	\end{itemize}
	\textbf{Recordkeeping Systems Requirements:}
	\begin{itemize}
		\item Secure access (limited to need-to-know)
		\item Audit trail of who accessed SAR files
		\item Tamper-proof records
		\item Ability to produce records upon request
		\item Backup and disaster recovery
	\end{itemize}
	\redflag{Records cannot be deleted until retention period expires, regardless of customer request.}
\end{frame}

% ============================================================
% SECTION 7: RESPONDING TO LAW ENFORCEMENT (6 SLIDES)
% ============================================================

\begin{frame}{6.7: Responding to Subpoenas and LE Requests}
	\begin{block}{Core Concept}
		Financial institutions have a legal obligation to respond to valid law enforcement requests.
	\end{block}
	\textbf{Key Principles:}
	\begin{itemize}
		\item Comply promptly within specified timeframe
		\item Maintain internal record of everything produced
		\item Do NOT notify customer unless subpoena requires disclosure
		\item Seek legal advice if uncertain about scope or validity
		\item Do not destroy records after receiving subpoena (legal hold)
		\item Produce only what subpoena requests (not everything)
		\item Redact privileged or unrelated information
	\end{itemize}
	\redflag{What NOT to do: Destroy records, notify customer to warn them, refuse to comply without legal basis.}
\end{frame}

\begin{frame}{6.7.1: Subpoena with Non-Disclosure Provision}
	\scenario{Bank Receives Grand Jury Subpoena}
	Bank receives federal grand jury subpoena demanding all account records, wire logs, and signature cards for corporate customer under investigation. Subpoena includes non-disclosure provision prohibiting bank from notifying customer.
	
	\vspace{0.1cm}
	\textbf{How to Analyze:}
	\begin{enumerate}
		\item Bank must comply with subpoena
		\item Non-disclosure provision means bank CANNOT tell customer
		\item Bank must maintain internal record of everything produced
		\item Bank should not destroy any responsive records
		\item Legal hold must be implemented
	\end{enumerate}
	\trap{Bank should notify customer so they can respond to subpoena.}
	\correct{Comply promptly, maintain internal records, do not violate non-disclosure provision.}
\end{frame}

\begin{frame}{6.7.2: Keep Open Request - Balancing Risks}
	\scenario{Law Enforcement Asks Bank to Keep Account Open}
	Bank filed three SARs on account over six months and decided to close relationship. Before closing, federal agent contacts bank in writing, requesting bank keep account open for additional 90 days for ongoing undercover investigation. Agent says closing would alert subjects.
	
	\vspace{0.1cm}
	\textbf{How to Analyze:}
	\begin{enumerate}
		\item Bank should generally comply with written LE requests
		\item Request must be documented in writing
		\item Senior management approval should be obtained
		\item Bank should continue enhanced monitoring during extension
		\item Document decision to keep open or close
	\end{enumerate}
	\trap{Bank must close account immediately regardless of LE request.}
	\correct{Document request in writing, obtain senior management approval, keep open as requested with enhanced monitoring.}
\end{frame}

\begin{frame}{6.7.3: Legal Hold Requirements After Subpoena}
	\begin{block}{Legal Hold - What Must Be Preserved}
		Once a subpoena is received, a legal hold suspends normal record destruction.
	\end{block}
	\textbf{Legal Hold Process:}
	\begin{itemize}
		\item Identify all relevant records (customer files, transaction data, communications)
		\item Notify all relevant departments (compliance, IT, records management)
		\item Suspend any scheduled destruction of records
		\item Preserve records in original format
		\item Maintain chain of custody
		\item Document legal hold implementation
		\item Release hold only after authorized by legal counsel
	\end{itemize}
	\textbf{Common Legal Hold Failures:}
	\begin{itemize}
		\item Automated record destruction continues after subpoena
		\item Relevant departments not notified
		\item No documentation of hold implementation
		\item Premature release of hold
		\item Failure to preserve electronic records (emails, logs)
	\end{itemize}
	\redflag{Destroying records after receiving subpoena can result in spoliation sanctions.}
\end{frame}

% ============================================================
% SECTION 8: DE-RISKING AND FINANCIAL INCLUSION (6 SLIDES)
% ============================================================

\begin{frame}{6.8: De-Risking - What It Is and Why Problematic}
	\begin{block}{Core Concept}
		De-risking = indiscriminate termination of entire categories of customer relationships without individual risk assessment.
	\end{block}
	\textbf{Why De-Risking is Problematic:}
	\begin{itemize}
		\item Drives legitimate activity into unregulated, informal channels
		\item Harms financial inclusion (remittance corridors, developing economies)
		\item Violates risk-based approach (requires individual assessment)
		\item May contradict regulatory expectations
		\item Can concentrate risk rather than reduce it
		\item May violate anti-discrimination laws in some jurisdictions
	\end{itemize}
	\textbf{Key Distinction:}
	\begin{itemize}
		\item De-risking (bad): Blanket termination by geography, customer type, or category without individual assessment
		\item Justified exit (good): Individual, documented termination based on specific risk assessment of that relationship
	\end{itemize}
	\redflag{FATF and Wolfsberg guidance explicitly discourages blanket de-risking.}
\end{frame}

\begin{frame}{6.8.1: Blanket Termination vs. Individual Assessment}
	\scenario{Bank Terminates All African Correspondent Relationships}
	Global bank terminates all correspondent relationships with banks headquartered in Sub-Saharan Africa. Termination applies to all 22 respondent banks, including three that passed recent FATF mutual evaluations with Compliant ratings. Bank conducted no individual risk assessments. Compliance director states: The region is high-risk, so we are exiting entirely.
	
	\vspace{0.1cm}
	\textbf{How to Analyze:}
	\begin{enumerate}
		\item Blanket termination without individual assessment violates RBA
		\item Compliant banks terminated along with non-compliant ones
		\item Drives legitimate activity out of formal system
		\item FATF and Wolfsberg guidance discourages this
		\item May harm legitimate remittance and trade flows
	\end{enumerate}
	\trap{Banks have absolute discretion to choose correspondent partners.}
	\correct{RBA requires individual, documented risk assessments. Blanket termination without assessment is de-risking and discouraged.}
\end{frame}

\begin{frame}{6.8.2: Documented Risk-Based Exit - Permissible}
	\scenario{Bank Terminates Specific Respondent After Assessment}
	Bank terminates correspondent relationship with specific respondent after: (1) individual risk assessment; (2) identifying respondent services 14 unlicensed MSBs; (3) issuing two unanswered RFIs; (4) documenting risk assessment and decision. Other banks in same jurisdiction retained.
	
	\vspace{0.1cm}
	\textbf{How to Analyze:}
	\begin{enumerate}
		\item Individual assessment conducted (not blanket termination)
		\item Specific risk factors identified
		\item Attempts to remediate (RFIs) were made
		\item Decision documented
		\item Other banks in same jurisdiction retained
	\end{enumerate}
	\trap{Any termination of correspondent relationship is de-risking.}
	\correct{Justified risk-based exit based on individual documented assessment is permissible and NOT de-risking.}
\end{frame}

\begin{frame}{6.8.3: FATF Guidance on De-Risking}
	\begin{block}{FATF's Position on De-Risking}
		FATF has issued specific guidance addressing the de-risking phenomenon.
	\end{block}
	\textbf{FATF Guidance Key Points:}
	\begin{itemize}
		\item Broad, categorical termination of entire classes of customers is NOT the preferred risk management strategy under RBA
		\item Financial institutions should conduct individual, documented risk assessments
		\item Blanket exclusions based solely on geography or customer type are discouraged
		\item De-risking can drive activity into unregulated channels, making monitoring harder
		\item Financial inclusion is a legitimate consideration
	\end{itemize}
	\textbf{Wolfsberg Group Guidance:}
	\begin{itemize}
		\item Banks should avoid indiscriminate de-risking
		\item Where relationship terminated, decision should be documented and risk-based
		\item Not solely based on commercial considerations
		\item Attempts to remediate should be made before termination
	\end{itemize}
	\textbf{Regulatory Expectations:}
	\begin{itemize}
		\item Terminations will be reviewed for pattern of de-risking
		\item Banks expected to have documented rationale
		\item Regulators may question blanket terminations
	\end{itemize}
\end{frame}

% ============================================================
% SECTION 9: REQUESTS FOR INFORMATION (RFIs) (4 SLIDES)
% ============================================================

\begin{frame}{6.9: Requests for Information (RFIs)}
	\begin{block}{Core Concept}
		RFIs are formal requests sent to customers or respondent banks asking for clarification or documentation of suspicious activity.
	\end{block}
	\textbf{Best Practices for RFIs:}
	\begin{itemize}
		\item Use when you need more information to resolve an alert
		\item Document the RFI and any response (or lack thereof)
		\item Do NOT disclose that an investigation or SAR is pending (tipping-off)
		\item Failure to respond is itself a red flag
		\item Set reasonable deadline for response (typically 5-10 business days)
		\item Escalate if no response received
		\item Maintain copy of RFI in case file
	\end{itemize}
	\redflag{Critical: RFIs must be phrased neutrally. Do NOT say we are investigating you for money laundering.}
\end{frame}

\begin{frame}{6.9.1: Neutral RFI vs. Tipping-Off Violation}
	\scenario{Analyst Tells Customer About AML Flag}
	Compliance analyst calls customer: Our AML system flagged your deposits as potential structuring. We need to know where you are getting all this cash, or we might have to file a report. Customer becomes angry and hangs up.
	
	\vspace{0.1cm}
	\textbf{How to Analyze:}
	\begin{enumerate}
		\item Mentioning AML system and potential structuring is tipping-off
		\item Mentioning file a report is explicit tipping-off
		\item Customer now knows they are under investigation
		\item RFI was phrased improperly
	\end{enumerate}
	\trap{Analyst was honest with customer, which is good customer service.}
	\correct{RFIs must be neutral. Example: We noticed increase in cash deposits. Can you provide documentation explaining source of funds? No mention of AML or reporting.}
\end{frame}

\begin{frame}{6.9.2: Proper RFI Wording and Documentation}
	\begin{block}{Examples of Proper RFI Wording}
		Neutral language that does not tip off the customer.
	\end{block}
	\textbf{Proper Wording Examples:}
	\begin{itemize}
		\item We noticed an increase in wire transfers from your account. Please provide documentation explaining the business purpose of these transactions.
		\item Please confirm the nature of your relationship with [counterparty name] and provide any supporting agreements or invoices.
		\item Our records show cash deposits totaling [amount] over the past [period]. Please provide documentation supporting the source of these funds.
		\item Please update your customer profile to reflect recent changes in your business activities.
	\end{itemize}
	\textbf{Documentation Requirements:}
	\begin{itemize}
		\item Copy of RFI sent to customer
		\item Date RFI was sent
		\item Method of delivery (email, letter, phone log)
		\item Any response received (or note that no response received)
		\item Follow-up actions taken
		\item Escalation if no response
	\end{itemize}
	\textbf{When to Escalate Without RFI:}
	\begin{itemize}
		\item Sending RFI would tip off customer (suspicion is very strong)
		\item Customer has history of destroying records when asked
		\item Law enforcement requests no contact
		\item Immediate SAR filing is more appropriate
	\end{itemize}
\end{frame}

% ============================================================
% SECTION 10: AML TRAINING REQUIREMENTS (6 SLIDES)
% ============================================================

\begin{frame}{6.10: AML Training - Role-Based Requirements}
	\begin{block}{Core Concept}
		Training must be tailored to the specific roles and risks of each employee group.
	\end{block}
	\textbf{Role-Based Training Requirements:}
	\begin{itemize}
		\item Tellers: Cash structuring, CTR triggers, identifying suspicious customer behavior, red flags
		\item Relationship Managers: PEP identification, EDD requirements, red flags in customer profiles, source of wealth
		\item Board of Directors: Governance oversight, risk strategy, regulatory expectations, enforcement trends
		\item IT/Technology: Data security, system controls, audit trail requirements, data quality
		\item Compliance Staff: Advanced investigation techniques, SAR writing, regulatory updates, emerging typologies
		\item Senior Management: Program oversight, resource allocation, escalation procedures, regulatory risk
		\item New Hires: Initial training before any customer contact
	\end{itemize}
	\redflag{Deficiency: The same training module delivered to all employees regardless of role.}
\end{frame}

\begin{frame}{6.10.1: One-Size-Fits-All Training Failure}
	\scenario{Bank Delivers Same Training to All 5,000 Employees}
	Bank delivers same two-hour online AML training module to all 5,000 employees: tellers, RMs, Board, IT staff. Training covers technical details of transaction monitoring thresholds. External auditor rates program ineffective because Board does not need threshold calculations, tellers received no cash structuring training.
	
	\vspace{0.1cm}
	\textbf{How to Analyze:}
	\begin{enumerate}
		\item Tellers need training on cash structuring and CTRs, not model calibration
		\item Board members need governance training, not technical threshold details
		\item One-size-fits-all training is ineffective
		\item Bank must develop role-specific modules
	\end{enumerate}
	\trap{All employees need same training to ensure consistent institutional coverage.}
	\correct{Training must be tailored to specific roles, responsibilities, and ML/TF risks. Regulators expect role-based training.}
\end{frame}

\begin{frame}{6.10.2: Board Training - Governance Not Operations}
	\scenario{Board Training Covers CTR Filing and Alert Review}
	Bank's annual Board AML training covers: (1) how to identify suspicious transactions in alerts; (2) how to file Currency Transaction Reports; (3) AML risk appetite and governance oversight; (4) regulatory expectations for Board oversight.
	
	\vspace{0.1cm}
	\textbf{How to Analyze:}
	\begin{enumerate}
		\item Board members do NOT file CTRs or review alerts (operational)
		\item Board training should cover risk appetite, governance, oversight
		\item Items 1 and 2 are inappropriate for Board training
		\item Items 3 and 4 are appropriate for Board training
	\end{enumerate}
	\trap{Board members need to understand all aspects of AML program, including operations.}
	\correct{Board training should focus on governance, oversight, risk strategy, and regulatory expectations, not operational tasks.}
\end{frame}

\begin{frame}{6.10.3: Training Frequency and Documentation}
	\begin{block}{Training Frequency Requirements}
		Initial training AND periodic refresher training are required.
	\end{block}
	\textbf{Frequency Requirements:}
	\begin{itemize}
		\item New hire training: Before any customer contact or within 30 days of hire
		\item Annual refresher training: At least once per calendar year
		\item Role change training: When employee moves to new position with different risks
		\item Policy change training: When material changes to AML program occur
		\item Remedial training: For employees who fail testing or make errors
		\item Regulatory update training: As needed when new regulations issued
	\end{itemize}
	\textbf{Documentation Requirements:}
	\begin{itemize}
		\item Training attendance records (who attended)
		\item Training content (what was covered, slides, materials)
		\item Training dates and duration
		\item Testing results (if applicable)
		\item Certificates of completion
		\item Records of remedial training
		\item Retention: Minimum 5 years
	\end{itemize}
	\textbf{Training Effectiveness Assessment:}
	\begin{itemize}
		\item Pre- and post-training testing
		\item Scenario-based testing
		\item Observations of employee behavior
		\item Audit of training program
		\item Tracking of repeated errors
	\end{itemize}
	\redflag{Training without effectiveness assessment may not satisfy regulatory requirements.}
\end{frame}

% ============================================================
% SECTION 11: WHISTLEBLOWER PROTECTIONS (4 SLIDES)
% ============================================================

\begin{frame}{6.11: Whistleblower Protections}
	\begin{block}{Core Concept}
		Laws protect employees who report AML violations from retaliation.
	\end{block}
	\textbf{Key Protections (United States):}
	\begin{itemize}
		\item Bank Secrecy Act (31 U.S.C. Section 5328) prohibits retaliation for reporting BSA violations
		\item Dodd-Frank Act Section 922 (SEC-related disclosures)
		\item FIRREA (financial institution anti-retaliation)
		\item Sarbanes-Oxley Act (public company whistleblowers)
		\item State law protections (varies by jurisdiction)
	\end{itemize}
	\textbf{What Constitutes Retaliation:}
	\begin{itemize}
		\item Termination or firing
		\item Demotion or reduction in responsibilities
		\item Suspension without pay
		\item Harassment or hostile work environment
		\item Negative performance reviews (PIPs) without prior issues
		\item Exclusion from meetings or opportunities
		\item Transfer to undesirable location
		\item Blacklisting within industry
	\end{itemize}
	\redflag{Retaliation can result in significant damages, reinstatement, and legal fees for bank.}
\end{frame}

\begin{frame}{6.11.1: Retaliatory PIP - Complete Analysis}
	\scenario{Analyst Receives PIP After Internal Report}
	Compliance analyst reports internally that senior management is suppressing SAR filings on high-revenue corporate client. Three weeks later, analyst receives performance improvement plan citing communication issues - her first negative review in eight years. She believes PIP is retaliatory.
	
	\vspace{0.1cm}
	\textbf{How to Analyze:}
	\begin{enumerate}
		\item Timing (3 weeks after report) suggests retaliation
		\item First negative review in 8 years is suspicious
		\item BSA prohibits retaliation for reporting BSA violations
		\item She can file complaint with DOL/OSHA within 180 days
		\item Internal reports are protected, not just external
	\end{enumerate}
	\trap{Internal reports are not protected; only external reports to FinCEN are protected.}
	\correct{BSA prohibits retaliation for internal reports as well. Document everything and file complaint with DOL/OSHA.}
\end{frame}

\begin{frame}{6.11.2: How to File a Whistleblower Complaint}
	\begin{block}{Whistleblower Complaint Process}
		Employees have specific channels to file retaliation complaints.
	\end{block}
	\textbf{Filing a Complaint (United States):}
	\begin{itemize}
		\item File with Department of Labor, Occupational Safety and Health Administration (OSHA)
		\item Deadline: Within 180 days of the adverse employment action
		\item No filing fee
		\item Can file anonymously or with attorney representation
		\item OSHA investigates and can order reinstatement, back pay, and damages
		\item If OSHA does not act, employee can take case to federal court
	\end{itemize}
	\textbf{What to Document:}
	\begin{itemize}
		\item Copy of original report (with date)
		\item Any evidence of retaliation (PIP, emails, witness statements)
		\item Performance reviews before and after report
		\item Timeline of events
		\item Any communications with management about the report
	\end{itemize}
	\textbf{Protection from Retaliation:}
	\begin{itemize}
		\item Bank cannot retaliate for filing complaint
		\item Additional remedies available if bank retaliates against whistleblower complaint
		\item Whistleblower can recover legal fees if successful
	\end{itemize}
	\redflag{Employers cannot require employees to sign confidentiality agreements waiving whistleblower rights.}
\end{frame}

% ============================================================
% SECTION 12: SUPPORTING DEPARTMENTS (6 SLIDES)
% ============================================================

\begin{frame}{6.12: Supporting Departments - Data Privacy and Cybersecurity}
	\begin{block}{Core Concept}
		AML compliance does not operate in isolation. Other departments enable effective AFC programs.
	\end{block}
	\textbf{Key Supporting Functions and Their Roles:}
	\begin{itemize}
		\item Data Privacy: Ensures customer data protected while enabling lawful information sharing
		\item Cybersecurity: Protects AML systems from attack; identifies cyber-related red flags (ransomware, unauthorized access)
		\item IT/Technology: Maintains transaction monitoring systems, data feeds, audit trails, system availability
		\item Human Resources: Employee due diligence, insider threat monitoring, background checks
		\item Vendor Management: Third-party due diligence for outsourced AML functions
		\item Legal: Advice on regulatory requirements, subpoena response, legal hold management
		\item Internal Audit: Independent testing of AML controls (Third Line)
		\item Risk Management: Enterprise risk assessment integration
	\end{itemize}
	\redflag{AML cannot function effectively without strong support from these departments.}
\end{frame}

\begin{frame}{6.12.1: Data Privacy vs. AML - GDPR Exception}
	\scenario{Bank Receives GDPR Right to Erasure Request}
	European bank receives GDPR right to erasure request from customer demanding deletion of all personal data, including KYC documents and transaction records. Customer was subject of SAR filed 18 months ago. Data Privacy Officer argues GDPR requires deletion. AML Officer argues AML recordkeeping requires retention.
	
	\vspace{0.1cm}
	\textbf{How to Analyze:}
	\begin{enumerate}
		\item GDPR Article 17(3) exempts data retention required by law
		\item AML recordkeeping (minimum 5 years) is legal obligation
		\item Bank can lawfully refuse to delete AML-required records
		\item SAR confidentiality also prohibits acknowledging SAR to customer
	\end{enumerate}
	\trap{GDPR right to erasure is absolute and overrides all other laws.}
	\correct{AML retention requirements are legal exception to GDPR erasure. Bank may refuse deletion for records required to retain.}
\end{frame}

\begin{frame}{6.12.2: Cybersecurity and AML - Ransomware Red Flag}
	\scenario{Customer Pays Ransom to SDN-Listed Wallet}
	Cybersecurity team detects corporate customer's systems compromised by ransomware. Customer pays ransom in Bitcoin to wallet address appearing on OFAC SDN list. Cybersecurity team notifies AML team.
	
	\vspace{0.1cm}
	\textbf{How to Analyze:}
	\begin{enumerate}
		\item Paying ransom to SDN-listed wallet is OFAC violation
		\item Transaction may also constitute money laundering
		\item Bank should file SAR
		\item Bank should consider whether to block transaction
		\item AML and cybersecurity must coordinate
	\end{enumerate}
	\trap{Bank should process payment to avoid harm to customer.}
	\correct{AML and cybersecurity teams must coordinate. Paying ransom to SDN-listed wallet requires sanctions review and potential SAR.}
\end{frame}

\begin{frame}{6.12.3: IT/Technology - System Availability and Integrity}
	\begin{block}{IT's Role in AML Compliance}
		IT maintains the systems that enable AML monitoring. Failures become compliance failures.
	\end{block}
	\textbf{IT Responsibilities for AML:}
	\begin{itemize}
		\item System availability (monitoring must run 24/7/365)
		\item Data integrity (no corruption, accurate processing)
		\item Audit trails (who accessed what, when, from where)
		\item Change management (documented approvals for system changes)
		\item Disaster recovery (backup systems, recovery time objectives)
		\item Security (access controls, preventing unauthorized changes)
		\item Performance (system processes alerts within required timeframes)
	\end{itemize}
	\textbf{Common IT Failures with AML Impact:}
	\begin{itemize}
		\item System outages during critical monitoring windows
		\item Data feed failures without alerting
		\item Unauthorized changes to thresholds or scenarios
		\item Loss of audit trail data
		\item Slow performance causing alert backlog
		\item Security breach allowing manipulation
	\end{itemize}
	\redflag{IT failures are compliance failures when they impact AML monitoring effectiveness.}
\end{frame}

% ============================================================
% SECTION 13: EMPLOYEE DUE DILIGENCE AND INSIDER THREATS (4 SLIDES)
% ============================================================

\begin{frame}{6.13: Employee Due Diligence and Insider Threats}
	\begin{block}{Core Concept}
		Insider threats (employees misusing their access) are a significant AML risk.
	\end{block}
	\textbf{Key Controls for Insider Threat Management:}
	\begin{itemize}
		\item Background checks on hiring (criminal, credit, employment verification)
		\item Ongoing employee monitoring (not just customer monitoring)
		\item Segregation of duties (no single employee controls end-to-end)
		\item Access controls and audit trails (who accessed what)
		\item Mandatory vacation policies (detect ongoing fraud)
		\item Job rotation for sensitive positions
		\item Employee hotline for reporting suspicious colleagues
		\item Monitoring for lifestyle changes (where permitted by law)
	\end{itemize}
	\redflag{Red flags: Employee living beyond means, refusing vacation, accessing accounts outside job scope, multiple overrides.}
\end{frame}

\begin{frame}{6.13.1: Insider Threat - Employee Override Analysis}
	\scenario{RM Overrides 23 Alerts for Same Client}
	AML analyst notices senior RM has overridden 23 transaction monitoring alerts in past year for same corporate client. Each override note states: Client is legitimate. RM has override access without supervisory approval. RM recently purchased new luxury car and vacation home.
	
	\vspace{0.1cm}
	\textbf{How to Analyze:}
	\begin{enumerate}
		\item 23 overrides for same client is unusual
		\item No supervisory approval for overrides is control weakness
		\item Lifestyle changes may indicate illicit enrichment
		\item Bank should investigate RM
		\item Overrides should require supervisor approval
		\item Pattern suggests RM may be compromised
	\end{enumerate}
	\trap{RM has authority to override alerts; no violation.}
	\correct{Multiple overrides for same client plus lifestyle changes require investigation. Overrides should require supervisory approval.}
\end{frame}

\begin{frame}{6.13.2: Mandatory Vacation Policy as Control}
	\begin{block}{Why Mandatory Vacation Detects Fraud}
		Employees committing fraud often avoid taking vacation to prevent detection.
	\end{block}
	\textbf{How Mandatory Vacation Works as Control:}
	\begin{itemize}
		\item Employee must take consecutive days off (typically 5-10 business days)
		\item Another employee performs their duties during absence
		\item Substitute may discover irregularities that employee was hiding
		\item Employee cannot cover up ongoing fraud while absent
		\item Common in financial institutions for high-risk roles
	\end{itemize}
	\textbf{Roles Where Mandatory Vacation is Critical:}
	\begin{itemize}
		\item Wire transfer processors
		\item Account opening staff
		\item Transaction monitoring override approvers
		\item Suspense account managers
		\item Reconciliation staff
		\item Customer data maintenance
	\end{itemize}
	\textbf{Red Flags Related to Vacation:}
	\begin{itemize}
		\item Employee never takes vacation
		\item Employee cancels vacation at last minute repeatedly
		\item Employee works during vacation
		\item Employee refuses to provide passwords for coverage
		\item Employee becomes anxious when vacation discussed
	\end{itemize}
	\redflag{Employees who never take vacation should be subject to enhanced scrutiny.}
\end{frame}

% ============================================================
% SECTION 14: FRONT OFFICE INTERACTION (4 SLIDES)
% ============================================================

\begin{frame}{6.14: Front Office Interaction - RMs and Compliance}
	\begin{block}{Core Concept}
		Relationship managers (First Line) and compliance (Second Line) have complementary but distinct roles.
	\end{block}
	\textbf{RM Responsibilities (First Line):}
	\begin{itemize}
		\item Know your customer (KYC) at onboarding
		\item Identify and report unusual activity to compliance
		\item Respond to RFIs from compliance
		\item Update customer information when changes occur
		\item Provide context on customer's legitimate business
		\item NOT make SAR filing decisions
	\end{itemize}
	\textbf{Compliance Responsibilities (Second Line):}
	\begin{itemize}
		\item Investigate suspicious activity independently
		\item Make SAR filing decisions (RMs do NOT decide)
		\item Escalate material findings to management and Board
		\item Cannot be overridden by RM for commercial reasons
		\item Must maintain independence from revenue pressure
	\end{itemize}
	\redflag{Critical: RMs cannot veto or prevent SAR filings based on commercial importance of client.}
\end{frame}

\begin{frame}{6.14.1: RM Who Blocks SAR - Commercial Pressure}
	\scenario{RM Argues Against SAR Due to Client Revenue}
	Compliance analyst concludes reasonable suspicion of money laundering on high-revenue corporate client generating \$5 million annual fees. RM argues: This client is too important. We cannot file a SAR based on suspicion. Get more evidence first. CCO under pressure from RM delays filing.
	
	\vspace{0.1cm}
	\textbf{How to Analyze:}
	\begin{enumerate}
		\item Reasonable suspicion is legal standard for SAR
		\item Commercial importance is NOT valid reason to delay or avoid filing
		\item CCO has independent obligation to file
		\item Delaying filing is regulatory violation
		\item RM attempting to improperly influence compliance
	\end{enumerate}
	\trap{RM commercial concern is legitimate and should be considered.}
	\correct{Compliance decisions on SAR filing must be independent of commercial considerations. CCO must file upon reasonable suspicion regardless of client revenue.}
\end{frame}

\begin{frame}{6.14.2: Proper RM-Compliance Interaction}
	\begin{block}{Healthy RM-Compliance Relationship}
		When functioning properly, RM and compliance work together while maintaining appropriate separation.
	\end{block}
	\textbf{RM Should Provide:}
	\begin{itemize}
		\item Context about customer's legitimate business
		\item Explanation of unusual transactions if available
		\item Documentation supporting legitimate activity
		\item Updated customer information
		\item Early warning of potential red flags
	\end{itemize}
	\textbf{RM Should NOT Do:}
	\begin{itemize}
		\item Pressure compliance to close alerts
		\item Threaten to escalate to senior management to override compliance
		\item Withhold information that might lead to SAR
		\item Discuss SAR filing with customer
		\item Make independent determination of suspicious activity
	\end{itemize}
	\textbf{Compliance Should Provide:}
	\begin{itemize}
		\item Clear guidance on what constitutes suspicious activity
		\item Training on red flags
		\item Timely response to RM inquiries
		\item Feedback on SAR outcomes (where permitted)
		\item Respect for legitimate customer relationships
	\end{itemize}
	\redflag{If RM is pressuring compliance to avoid SAR filing, escalate to higher authority.}
\end{frame}

% ============================================================
% SECTION 15: THIRD-PARTY AND VENDOR DUE DILIGENCE (4 SLIDES)
% ============================================================

\begin{frame}{6.15: Third-Party and Vendor Due Diligence}
	\begin{block}{Core Concept}
		Banks remain responsible for outsourced AML functions. Vendor risk must be managed.
	\end{block}
	\textbf{Key Requirements for Vendor Management:}
	\begin{itemize}
		\item Due diligence on vendors before engagement
		\item Contractual obligations for AML compliance
		\item Right to audit vendor performance
		\item Ongoing monitoring of vendor effectiveness
		\item Bank retains ultimate regulatory responsibility
		\item Cannot outsource accountability
	\end{itemize}
	\textbf{Vendors Requiring AML Due Diligence:}
	\begin{itemize}
		\item Transaction monitoring system providers
		\item Sanctions screening vendors
		\item KYC/CIP verification services
		\item Customer identification document vendors
		\item Case management system providers
		\item Training providers
		\item Independent testing firms
	\end{itemize}
	\redflag{Cannot outsource responsibility: Even if vendor performs AML functions, bank is still accountable to regulators.}
\end{frame}

\begin{frame}{6.15.1: Outsourced Monitoring Failure - Liability}
	\scenario{Vendor System Fails for 8 Months}
	Bank outsources transaction monitoring to third-party vendor. Vendor system fails to generate alerts for 8 months due to technical error. Bank's internal compliance team did not perform any validation of vendor outputs. When discovered, vendor claims liability. Bank argues vendor is responsible.
	
	\vspace{0.1cm}
	\textbf{How to Analyze:}
	\begin{enumerate}
		\item Bank is ultimately responsible to regulators
		\item Bank cannot delegate liability to vendor
		\item Bank should have performed validation and reconciliation
		\item Bank must remediate (review missed period, file SARs)
		\item Bank must address vendor oversight failure
	\end{enumerate}
	\trap{Vendor contract makes them liable for monitoring failures.}
	\correct{Bank retains regulatory responsibility regardless of vendor contracts. Bank must have vendor oversight and validation controls.}
\end{frame}

\begin{frame}{6.15.2: Vendor Due Diligence Requirements}
	\begin{block}{Pre-Engagement and Ongoing Vendor Diligence}
		Comprehensive due diligence required before and during vendor relationship.
	\end{block}
	\textbf{Pre-Engagement Due Diligence:}
	\begin{itemize}
		\item Vendor background and reputation
		\item Financial stability of vendor
		\item Vendor's own AML program
		\item Vendor's data security practices
		\item Vendor's subcontractor relationships
		\item References from other bank clients
		\item Regulatory history (any enforcement actions)
		\item Business continuity and disaster recovery
	\end{itemize}
	\textbf{Contract Requirements:}
	\begin{itemize}
		\item Right to audit vendor (on-site and remote)
		\item Service level agreements (uptime, accuracy)
		\item Data ownership and return provisions
		\item Breach notification requirements
		\item Subcontractor approval rights
		\item Liability and indemnification provisions
		\item Termination rights for non-performance
		\item Confidentiality and data protection
	\end{itemize}
	\textbf{Ongoing Monitoring:}
	\begin{itemize}
		\item Regular performance reviews
		\item Service level monitoring
		\item Annual vendor risk assessment
		\item Periodic on-site audits
		\item Vendor security assessment updates
	\end{itemize}
\end{frame}

% ============================================================
% SECTION 16: ESCALATION AND OFFBOARDING (4 SLIDES)
% ============================================================

\begin{frame}{6.16: Escalation and Offboarding}
	\begin{block}{Core Concept}
		When a customer poses unacceptable risk, bank must escalate through governance and potentially offboard.
	\end{block}
	\textbf{Escalation Path for High-Risk Customers:}
	\begin{enumerate}
		\item Analyst identifies suspicious activity
		\item Investigation conducted and documented
		\item If SAR filed, notify AML committee
		\item If multiple SARs or high risk, escalate to senior management
		\item Board notification for material risks
		\item Decision to offboard requires senior management approval
		\item Offboarding executed according to procedures
	\end{enumerate}
	\textbf{Documentation Required for Escalation:}
	\begin{itemize}
		\item Summary of suspicious activity
		\item Investigation findings
		\item SARs filed (if any)
		\item Risk assessment of relationship
		\item Attempts to remediate (RFIs, enhanced monitoring)
		\item Recommendation (close or continue)
		\item Approval documentation
	\end{itemize}
	\redflag{Offboarding must be documented and risk-based, not retaliatory or discriminatory.}
\end{frame}

\begin{frame}{6.16.1: Offboarding Decision After Multiple SARs}
	\scenario{Bank Has Filed 7 SARs on Account Over 14 Months}
	Bank filed 7 SARs on account over 14 months. Account shows clear patterns of structured deposits followed by wires to offshore shell companies. Compliance committee recommends offboarding. RM argues: No conviction has occurred. Client is innocent until proven guilty. We should keep the account.
	
	\vspace{0.1cm}
	\textbf{How to Analyze:}
	\begin{enumerate}
		\item Banks are not courts of law
		\item Standard for offboarding is risk, not conviction
		\item 7 SARs on one account is clear evidence of elevated risk
		\item Bank can (and often should) offboard
		\item RM argument confuses criminal standard with risk management
	\end{enumerate}
	\trap{Customer is innocent until proven guilty, so bank should keep account.}
	\correct{Offboarding is risk management decision, not criminal conviction. Multiple SARs justify offboarding regardless of conviction status.}
\end{frame}

\begin{frame}{6.16.2: Offboarding Procedures and Best Practices}
	\begin{block}{Proper Offboarding Execution}
		Offboarding must be done carefully to avoid legal and reputational risk.
	\end{block}
	\textbf{Offboarding Steps:}
	\begin{enumerate}
		\item Document decision to offboard with rationale
		\item Obtain required approvals (compliance committee, senior management)
		\item Notify law enforcement if ongoing investigation
		\item Provide notice to customer (standard account closure notice)
		\item Do NOT disclose reason if related to SAR (tipping-off)
		\item Return any remaining funds to customer
		\item File SAR if closure itself is suspicious
		\item Maintain records of offboarding decision
	\end{enumerate}
	\textbf{What NOT to Do When Offboarding:}
	\begin{itemize}
		\item Do not tell customer we are closing account because of SAR
		\item Do not close account to retaliate against whistleblower customer
		\item Do not close account without documentation
		\item Do not delay returning customer funds
		\item Do not offboard solely based on protected characteristics
	\end{itemize}
	\textbf{Alternative to Full Offboarding:}
	\begin{itemize}
		\item Restrict account activity (no wires, no new deposits)
		\item Enhanced monitoring
		\item Lower transaction limits
		\item Require prior approval for certain transactions
		\item Convert to cash management only
	\end{itemize}
	\redflag{If law enforcement requests account remain open, document and comply.}
\end{frame}


% ============================================================
% MISSING SLIDES FOR DOMAIN C (Building an AFC Program)
% Insert after the "Three Lines of Defense" section (around slide 6.1.10)
% ============================================================

\begin{frame}{6.1.11: AML Program Governance Documentation Requirements}
	\begin{block}{Written Policies, Procedures, and Controls}
		The BSA/AML compliance program must be reduced to writing and approved by the Board. Documentation is not optional.
	\end{block}
	\textbf{Mandatory Governance Documents:}
	\begin{itemize}
		\item Written AML/CFT Policy approved by Board (initial and all material changes)
		\item Written CIP (Customer Identification Program) procedures
		\item Written CDD/EDD procedures (including beneficial ownership)
		\item Written SAR filing procedures (internal decision-making, timeliness)
		\item Written independent testing procedures and scope
		\item Written employee training program (role-based curriculum)
		\item Written risk assessment methodology (Enterprise-Wide Risk Assessment)
		\item Written escalation procedures (compliance to management to Board)
		\item Written recordkeeping procedures (retention periods, secure storage)
	\end{itemize}
	\textbf{Documentation Best Practices:}
	\begin{itemize}
		\item Version control (date, author, approval)
		\item Board meeting minutes reflecting approval and discussion
		\item Centralized policy repository with access controls
		\item Regular policy review (at least annually)
		\item Tracking of employee acknowledgments
	\end{itemize}
	\redflag{Common exam trap: A verbal policy or unwritten procedure does NOT satisfy regulatory requirements.}
\end{frame}

\begin{frame}{6.1.12: Board Meeting Minutes - What Must Be Documented}
	\scenario{Board Verbally Approves AML Policy Without Minutes}
	Board chair verbally approves updated AML policy during meeting. Secretary does not record vote or discussion in minutes. CCO argues the policy is effective because Board approved it. Regulator requests minutes and finds no record of approval, discussion, or any AML oversight.
	
	\vspace{0.1cm}
	\textbf{How to Analyze:}
	\begin{enumerate}
		\item Unwritten approval is insufficient — approval must be documented in minutes
		\item Minutes must show evidence of meaningful oversight, not rubber-stamping
		\item Discussion of AML risks, resource allocation, and key metrics must appear
		\item Regulators will treat missing minutes as evidence of Board non-engagement
	\end{enumerate}
	\trap{Verbal approval is still approval; minutes are administrative formality.}
	\correct{Board minutes must document AML approval, discussion, and ongoing oversight.}
\end{frame}

\begin{frame}{6.1.13: Chief Compliance Officer (CCO) / BSA Officer - Independence and Authority}
	\begin{block}{CCO Reporting Line and Authority}
		The CCO/BSA Officer must have sufficient authority, independence, and resources to perform their duties.
	\end{block}
	\textbf{Required Attributes of CCO Role:}
	\begin{itemize}
		\item Direct reporting access to Board or Board Committee (not filtered by business lines)
		\item Authority to allocate resources (staff, technology, training)
		\item Authority to override business decisions regarding customer relationships
		\item No reporting line to revenue-generating functions (e.g., Head of Sales, CEO as business line)
		\item Sufficient seniority to influence decisions
		\item Protection from retaliation for exercising duties
	\end{itemize}
	\textbf{Red Flags for CCO Independence Violation:}
	\begin{itemize}
		\item CCO reports to Head of Retail Banking (revenue function)
		\item Business lines control CCO's budget
		\item CCO is removed after filing SAR on high-revenue client
		\item CCO position is vacant for extended period
		\item Same person serves as CCO and Head of Sales
	\end{itemize}
	\redflag{If CCO can be overridden by business lines for commercial reasons, CCO lacks required authority.}
\end{frame}

\begin{frame}{6.1.14: CCO Reporting Line Violation - Reporting to Revenue Function}
	\scenario{CCO Reports to Head of International Banking (Revenue)}
	Bank's CCO reporting line is to Head of International Banking, who is compensated based on revenue growth. Head of International Banking instructs CCO to delay SAR filing on large corporate client because the client is negotiating a \$50 million credit facility. CCO's proposed budget was reduced after previous SAR filings affected revenue.
	
	\vspace{0.1cm}
	\textbf{How to Analyze:}
	\begin{enumerate}
		\item CCO must report to Board or independent committee, not revenue-generating management
		\item Revenue pressure on CCO compromises independence
		\item Delaying SAR for commercial reasons is a regulatory violation
		\item Budget reduction following SAR filings suggests retaliation
	\end{enumerate}
	\trap{Reporting to senior management is acceptable as long as CCO has Board access.}
	\correct{CCO must report to independent authority (Board/Audit Committee) without revenue conflicts. CCO cannot be subject to commercial override.}
\end{frame}

\begin{frame}{6.1.15: Program Resources - Staffing, Technology, and Budget}
	\begin{block}{Adequate Resources = Regulatory Requirement}
		Institutions must provide AML compliance with sufficient resources to fulfill its mandate.
	\end{block}
	\textbf{Resource Areas Requiring Adequacy:}
	\begin{itemize}
		\item Staffing: Sufficient analysts to review all alerts within required timeframes
		\item Technology: Monitoring systems, screening tools, case management, data infrastructure
		\item Training: Ongoing, role-based training for compliance and business staff
		\item External support: Legal, audit, third-party consultants as needed
		\item Budget: Competitive salaries to retain qualified talent
	\end{itemize}
	\textbf{Evidence of Resource Inadequacy (Exam Red Flags):}
	\begin{itemize}
		\item Persistent backlog of uninvestigated alerts (thousands aging >30 days)
		\item High analyst turnover (burnout, underpaid, under-supported)
		\item Manual processes for high-volume tasks that could be automated
		\item Delayed SAR filings due to insufficient reviewers
		\item Board rejects reasonable budget requests from CCO
		\item Staff not qualified for assigned duties (lack of training or experience)
	\end{itemize}
	\redflag{Under-resourcing is NOT a defense. Regulators will cite inadequate resources as a violation.}
\end{frame}

\begin{frame}{6.1.16: Regulatory Reporting Timelines - Key Deadlines}
	\begin{block}{Mandatory Reporting Deadlines (US Focus)}
		Different reports have different filing deadlines. Missing deadlines is a violation.
	\end{block}
	\textbf{Key Deadlines:}
	\begin{tabular}{p{0.35\textwidth}|p{0.30\textwidth}|p{0.30\textwidth}}
		\hline
		\rowcolor{lightgray}
		\textbf{Report Type} & \textbf{Filing Deadline} & \textbf{To Whom} \\
		\hline
		SAR (suspicious activity found) & Within 30 calendar days (60 days if no suspect identified) & FinCEN \\
		\hline
		CTR (cash >\$10,000 in one day) & Within 15 days & FinCEN \\
		\hline
		FBAR (foreign accounts >\$10,000 aggregate) & April 15 (with extension to October 15) & FinCEN \\
		\hline
		FinCEN 314(a) Response & Within 14 business days & FinCEN \\
		\hline
		Currency/Bearer Goods Report (CMIR) & At time of entry/exit & CBP \\
		\hline
	\end{tabular}
	\vspace{0.2cm}
	\redflag{SAR clock starts when suspicion is reached, not when investigation is complete.}
	\examfocus{Exam tests whether you know SAR deadline = 30 days (with 60-day extension for no suspect).}
\end{frame}

% ============================================================
% MISSING SLIDES FOR THIRD-PARTY RISK (Expanded)
% Insert after existing vendor slides (6.15.2)
% ============================================================

\begin{frame}{6.15.3: Correspondent Banking - Section 319(b) Record Production}
	\begin{block}{USA PATRIOT Act Section 319(b) - Foreign Bank Records}
		US correspondent banks must maintain records identifying agents for service of process for foreign respondent banks.
	\end{block}
	\textbf{Key Requirements of Section 319(b):}
	\begin{itemize}
		\item US banks must maintain records identifying the foreign bank's agent for service of legal process
		\item Agent must be located in the United States
		\item Records must be maintained for the duration of the correspondent relationship plus 5 years
		\item US government can demand records from foreign bank via this agent
		\item Failure to comply can result in termination of correspondent account
	\end{itemize}
	\textbf{Exam Scenario Facts:}
	\begin{itemize}
		\item US subpoena served on US bank requesting records of foreign respondent bank
		\item Foreign bank refuses to produce records
		\item US bank may be required to terminate the correspondent account
		\item Funds in the correspondent account may be subject to forfeiture
	\end{itemize}
	\redflag{Foreign bank's refusal to produce records under 319(b) is grounds for account termination.}
\end{frame}

\begin{frame}{6.15.4: Section 311 - Special Due Diligence for Primary Money Laundering Concerns}
	\begin{block}{USA PATRIOT Act Section 311 - Designated Institutions}
		FinCEN can designate foreign jurisdictions, institutions, or transaction types as "primary money laundering concerns."
	\end{block}
	\textbf{Section 311 Measures (Escalating Severity):}
	\begin{enumerate}
		\item Special Due Diligence: Enhanced recordkeeping, reporting, information collection
		\item Enhanced Due Diligence: Identify beneficial owners, monitor transactions, restrict accounts
		\item Prohibition on Accounts: U.S. banks cannot open or maintain correspondent accounts
	\end{enumerate}
	\textbf{Key Exam Points:}
	\begin{itemize}
		\item Section 311 targets foreign institutions (not U.S. banks directly)
		\item U.S. banks must comply with applicable 311 measures
		\item Designation is public and signals high risk
		\item Institutions may be removed from 311 list after remediation
	\end{itemize}
	\trap{U.S. banks can ignore Section 311 designations if they have no direct accounts with the named institution.}
	\correct{U.S. banks must apply special due diligence even for transactions that indirectly involve 311-designated entities.}
\end{frame}

\begin{frame}{6.15.5: Section 314(a) - Law Enforcement Information Requests}
	\begin{block}{FinCEN 314(a) - Mandatory Record Searches}
		FinCEN issues 314(a) requests on behalf of law enforcement. Financial institutions MUST search records and respond.
	\end{block}
	\textbf{314(a) Process:}
	\begin{enumerate}
		\item Law enforcement initiates request through FinCEN
		\item FinCEN issues 314(a) request to financial institutions
		\item Institutions search accounts and transactions for named subjects
		\item Responses due within 14 business days
		\item Confidentiality of request must be maintained
		\item Customer cannot be notified of the request
	\end{enumerate}
	\textbf{What Must Be Searched:}
	\begin{itemize}
		\item Current customer accounts
		\item Transactions within specified lookback period (typically 12 months)
		\item Accounts maintained for named individuals and entities
	\end{itemize}
	\textbf{What Happens with Positive Matches:}
	\begin{itemize}
		\item Institution reports match to FinCEN through secure system
		\item FinCEN forwards to requesting law enforcement agency
		\item Institution continues normal monitoring (SAR filing if activity suspicious)
	\end{itemize}
	\redflag{Failure to respond to 314(a) request within 14 business days is a regulatory violation.}
\end{frame}

\begin{frame}{6.15.6: Section 314(b) - Voluntary Information Sharing (Detailed)}
	\begin{block}{FinCEN 314(b) - The AML Information Sharing Safe Harbor}
		Financial institutions may share information with each other about suspected money laundering or terrorist financing.
	\end{block}
	\textbf{314(b) Requirements:}
	\begin{itemize}
		\item Institution must file notice with FinCEN to opt in (one-time)
		\item Sharing must be with other opted-in institutions
		\item Sharing must be for purpose of identifying and reporting ML/TF activity
		\item Institutions receive safe harbor from civil liability
		\item Sharing is NOT tipping-off
	\end{itemize}
	\textbf{What Can Be Shared:}
	\begin{itemize}
		\item Suspicious transaction patterns
		\item Identities of suspected persons or entities
		\item Account information (with appropriate safeguards)
		\item SARs (only with other opted-in institutions)
	\end{itemize}
	\textbf{What Cannot Be Shared:}
	\begin{itemize}
		\item Information for non-AML purposes (marketing, competitive intelligence)
		\item With institutions not opted into 314(b)
		\item With public or media
		\item With customers or subjects of investigation
	\end{itemize}
	\redflag{314(b) sharing is PERMISSIBLE and PROTECTED. Institutions cannot be penalized for good-faith sharing.}
\end{frame}

\begin{frame}{6.15.7: USA PATRIOT Act - Summary of Key Sections for CAMS}
	\begin{block}{Most Tested USA PATRIOT Act Sections in Domain C}
		Quick reference for exam preparation.
	\end{block}
	\begin{tabular}{p{0.15\textwidth}|p{0.25\textwidth}|p{0.55\textwidth}}
		\hline
		\rowcolor{lightgray}
		\textbf{Section} & \textbf{What It Does} & \textbf{Key Provision} \\
		\hline
		311 & Targets foreign ML concerns & Special due diligence; possible prohibition of accounts \\
		\hline
		312 & EDD for correspondent/private banking & Due diligence for foreign correspondent/PEP accounts \\
		\hline
		313 & Shell bank prohibition & Cannot maintain correspondent account for shell bank \\
		\hline
		314(a) & Law enforcement requests & Mandatory record search; 14-day response \\
		\hline
		314(b) & Voluntary information sharing & Safe harbor for sharing with other institutions \\
		\hline
		319(a) & Forfeiture of interbank accounts & Seizure of funds in correspondent accounts \\
		\hline
		319(b) & Record production & Foreign bank records via US agent \\
		\hline
		326 & CIP requirements & Customer identification program mandate \\
		\hline
		352 & AML program requirement & Written program, five pillars \\
		\hline
	\end{tabular}
	\examfocus{Know which section applies to which scenario. Section 311 = foreign institutions; 313 = shell banks; 314(a) = mandatory search; 314(b) = voluntary sharing.}
\end{frame}

% ============================================================
% MISSING SLIDES FOR RISK ASSESSMENT METHODOLOGY
% Insert after the Three Lines section (before Independent Testing)
% ============================================================

\begin{frame}{6.1.17: Enterprise-Wide Risk Assessment (EWRA) - Mandatory Components}
	\begin{block}{Risk Assessment as the Foundation}
		The AML program must be based on a documented, comprehensive risk assessment.
	\end{block}
	\textbf{Required Risk Factors (Must Address All):}
	\begin{multicols}{2}
		\begin{itemize}
			\item Products and services offered
			\item Customer types and segments
			\item Geographic locations (operations and customer base)
			\item Delivery channels (online, branch, mobile, agents)
			\item Transaction volumes and values
			\item Jurisdictions with weak AML regimes
			\item Politically Exposed Persons (PEPs)
			\item Non-profit organization exposure
			\item Correspondent banking relationships
			\item Private banking relationships
			\item Emerging technologies (crypto, etc.)
			\item Previous regulatory findings
		\end{itemize}
	\end{multicols}
	\textbf{Risk Assessment Outputs:}
	\begin{itemize}
		\item Inherent risk score (risk without controls)
		\item Control effectiveness rating
		\item Residual risk score (risk after controls)
		\item Action plan for residual risks exceeding tolerance
	\end{itemize}
	\redflag{EWRA must be UPDATED at least annually or after material changes. Static risk assessments are a deficiency.}
\end{frame}

\begin{frame}{6.1.18: Risk Assessment - Inherent vs. Residual Risk Calculation}
	\scenario{Bank Has High Inherent Risk But Strong Controls}
	Bank serves high-risk customers (MSBs, NPOs in conflict zones) and operates in high-risk jurisdictions (inherent risk = 8/10). Bank has strong controls: dedicated EDD team, transaction monitoring with low false negatives, monthly Board reporting, no regulatory findings (control effectiveness = 9/10). Residual risk = 3/10.
	
	\vspace{0.1cm}
	\textbf{How to Analyze:}
	\begin{enumerate}
		\item High inherent risk does NOT automatically mean deficient program
		\item Strong controls can reduce residual risk to acceptable levels
		\item Bank must document both inherent and residual risk
		\item Low residual risk with high inherent risk is acceptable if controls are truly effective
		\item Regulator will test whether controls are as strong as claimed
	\end{enumerate}
	\trap{High inherent risk is always unacceptable regardless of controls.}
	\correct{Risk-based approach accepts higher inherent risk if controls are proportionate and effective.}
\end{frame}

% ============================================================
% MISSING SLIDES FOR RECORDKEEPING (Expanded)
% Insert after Data Quality section
% ============================================================

\begin{frame}{6.4.5: Recordkeeping Requirements - Retention Periods}
	\begin{block}{Mandatory Record Retention (BSA/FATF)}
		Financial institutions must retain specified records for minimum periods.
	\end{block}
	\textbf{Core Retention Requirements:}
	\begin{itemize}
		\item Account opening records (CIP/KYC): 5 years after account closing
		\item Transaction records: 5 years from transaction date
		\item SAR filings and supporting documentation: 5 years from SAR filing date
		\item CTR filings: 5 years from filing date
		\item Funds transfer records (over \$3,000): 5 years
		\item Suspicious activity investigation records (even if no SAR filed): 5 years
		\item Training records: 5 years (track completions, content, dates)
		\item Independent testing reports: 5 years
		\item Board meeting minutes (AML-related): Permanent
	\end{itemize}
	\textbf{What "5 Years" Means:}
	\begin{itemize}
		\item Five years from the date of the last activity or filing, not from account opening
		\item For SARs: five years from SAR filing date
		\item Electronic records must be accessible and readable for full period
	\end{itemize}
	\redflag{Deleting records before retention period expires is a violation, even if customer requests deletion.}
\end{frame}

% ============================================================
% SUMMARY SLIDE FOR MISSING CONTENT
% Insert before or after existing conclusion slides
% ============================================================

\begin{frame}{Domain C - COMPLETE COVERAGE VERIFICATION}
	\begin{block}{All Missing Topics Now Added to Your Slides}
		The following slides have been added to complement your existing materials:
	\end{block}
	\begin{multicols}{2}
		\begin{itemize}
			\item Written policies and Board approval documentation
			\item Board meeting minutes requirements
			\item CCO independence and reporting line
			\item CCO authority over business lines
			\item Adequate resources (staffing, budget, technology)
			\item Regulatory reporting deadlines (SAR, CTR, FBAR)
			\item USA PATRIOT Act Section 319(b) (foreign bank records)
			\item USA PATRIOT Act Section 311 (primary ML concerns)
			\item USA PATRIOT Act Section 314(a) (LE information requests)
			\item USA PATRIOT Act Section 314(b) (voluntary sharing)
			\item Summary table of key USA PATRIOT sections
			\item Enterprise-Wide Risk Assessment (EWRA) components
			\item Inherent vs. residual risk calculation
			\item Recordkeeping retention periods (5-year rule)
		\end{itemize}
	\end{multicols}
	\vspace{0.2cm}
	\bluenote{These slides ensure full syllabus coverage for Domain C (30\% of CAMS exam).}
\end{frame}
% ============================================================
% CONCLUSION SLIDES (4 SLIDES)
% ============================================================

\begin{frame}{Domain 6 Summary: Top 30 Most Tested Concepts}
	\begin{multicols}{2}
		\begin{enumerate}
			\item Three Lines: First=Business, Second=Compliance, Third=Audit
			\item Independence: Audit cannot design controls it later tests
			\item Board approves program, CCO implements
			\item False negatives worse than false positives
			\item Over-filtering raises false negatives
			\item Model drift requires recalibration
			\item Silent data failure = monitoring gap
			\item Alert adjudication must be documented
			\item SAR standard = reasonable suspicion
			\item Tipping-off includes disclosing investigation
			\item 314(b) is exception to tipping-off
			\item Comply with subpoenas, do not notify customer
			\item Keep open requests: document, approve, comply
			\item De-risking = blanket termination (bad)
			\item Justified exit = individual assessment (good)
			\item RFIs must be neutral phrasing
			\item Training must be role-based
			\item Whistleblower protections exist for internal reports
			\item Vendor oversight: bank retains responsibility
			\item Offboarding is risk decision, not conviction
			\item Insider threats: monitor employees
			\item Mandatory vacation detects fraud
			\item RMs cannot block SARs for commercial reasons
			\item GDPR has AML retention exception
			\item Cybersecurity and AML coordinate on ransomware
			\item Data reconciliation prevents silent failures
			\item QA reviews ensure analyst consistency
			\item Legal hold required after subpoena
			\item FATF discourages blanket de-risking
			\item Board must be informed of material deficiencies
		\end{enumerate}
	\end{multicols}
\end{frame}

\begin{frame}{Critical Red Flags - Domain 6 Quick Reference}
	\begin{multicols}{2}
		\begin{itemize}
			\item Audit designing controls → Independence violation
			\item RM closing alerts → First Line overstepping
			\item No documentation in case files → Unverifiable
			\item 77 percent false negative rate → Material deficiency
			\item 34 percent missing data for 14 months → Data failure
			\item Teller says compliance reviewing → Tipping-off
			\item Blanket termination of African banks → De-risking
			\item RM blocks SAR for revenue → Improper pressure
			\item Same training for Board and tellers → Ineffective
			\item PIP 3 weeks after internal report → Retaliation
			\item Employee override without approval → Insider risk
			\item Offboarding delayed for more evidence → Misunderstanding
			\item CCO approves own policy → Board approval missing
			\item Management hides deficiency from Board → Violation
			\item Vendor failure with no oversight → Outsourcing failure
			\item Analyst assumes cash is normal → No investigation
			\item Different analysts different rates → No standards
			\item Thresholds raised 250 percent arbitrary → Over-filtering
			\item Never recalibrated model → Model drift
			\item No reconciliation controls → Data gap undetected
		\end{itemize}
	\end{multicols}
\end{frame}

\begin{frame}{Exam Day Reminders for Domain 6}
	\begin{itemize}
		\item Domain 6 is 30 percent of exam - second largest domain
		\item Three Lines of Defense is foundational - know who does what
		\item Independence is critical - audit cannot audit its own work
		\item SAR standard = reasonable suspicion, not proof
		\item Tipping-off includes disclosing investigation exists
		\item False negatives are worse than false positives
		\item De-risking = blanket termination (bad); justified exit = individual assessment (good)
		\item Training must be role-based
		\item Whistleblower protections exist for internal and external reports
		\item Banks cannot outsource regulatory responsibility to vendors
		\item Offboarding is risk decision, not conviction
		\item RMs cannot override compliance SAR decisions
		\item Board must approve AML program
		\item Material deficiencies must be reported to Board
		\item RFIs must be neutral - cannot mention AML or filing
		\item 314(b) sharing is permitted and not tipping-off
		\item Legal hold required after subpoena
		\item Data reconciliation prevents silent feed failures
		\item QA ensures analyst consistency
		\item Mandatory vacation detects insider fraud
	\end{itemize}
	\vspace{0.3cm}


\end{frame}

\end{document}